% SGA 3, Expose III, Loop 1 English reconstruction.
% Scope: 4.23 through the discussion following Corollary 4.30, stopping
% immediately before Corollary 4.31.
% Authority: Polo--Gille Exp3-14oct24.pdf, component-local pp. 67--70
% (combined-reader pp. 177--180; re-edition pp. 167--170; original
% side-pages 148--152). OCR is a locator only. English wording was compared
% with Jacob C. Reinhold's CC BY 4.0 Markdown lineage at commit
% e7a259f3f8608ad3edf9bf6eead3fd504dd2d23e. The Polo--Gille PDF controls
% prose, mathematics, numbering, editorial notes, and the two source-derived
% Loop-1 diagram images included below.

\setcounter{footnote}{120}

% Source page: Exp3 local p. 67; combined p. 177; re-edition p. 167;
% original side-page 148.
\medskip
\noindent\SGAInlineRefTarget{sga3:iii:paragraph:4-23}\textbf{4.23. ---}
仍在 \SGARefLink{sga3:iii:theorem:4-21}{4.21} 的假设下, 现在研究提升
\(Y_{\mathscr J}\) 的子群 \(Y\) 所成的集合在 \(X\) 的截面共轭作用下的性质.
若 \(x\) 是 \(X\) 在 \(S\) 上的截面, 且它诱导 \(X_{\mathscr J}\) 的单位截面,
则 \(x\) 所定义的内自同构 \(\operatorname{Int}(x)\) 将 \(X\) 中提升
\(Y_{\mathscr J}\) 的平坦子群映为 \(X\) 中提升 \(Y_{\mathscr J}\) 的平坦子群.
在 \SGARefLink{sga3:iii:theorem:4-21:item:ii}{4.21(ii)} 的条件下,
这些子群组成的集合是 \(Z^1(Y_0,N_0)\) 下的一个主齐性空间.
我们将看到, 此时 \(B^1(Y_0,N_0)\) 有一个子群 \(\Gamma\)
\footnote{编者注: 已将 \(Z^1(Y_0,N_0)\) 替换为 \(B^1(Y_0,N_0)\),
因为证明表明 \(\Gamma\) 是 \(B^1(Y_0,N_0)\) 的子群;
参见 \SGARefLink{sga3:iii:lemma:4-27}{4.27}--\SGARefLink{sga3:iii:corollary:4-29}{4.29}.},
使得 \(X\) 中两个在 \(S\) 上平坦且提升 \(Y_{\mathscr J}\) 的子群,
由一个诱导 \(X(S_{\mathscr J})\) 的单位元的 \(x\in X(S)\) 共轭,
当且仅当它们在 \(Z^1(Y_0,N_0)\) 中的\emph{差}属于 \(\Gamma\).
在最好的情形下, 我们将证明 \(\Gamma=B^1(Y_0,N_0)\).
因此, \(X\) 中提升 \(Y_{\mathscr J}\) 的平坦子群所成的集合,
模去诱导 \(X(S_{\mathscr J})\) 的单位元的元素 \(x\in X(S)\) 所给的共轭作用后,
或者为空, 或者是 \(H^1(Y_0,N_0)\) 下的主齐性空间
(参见 \SGARefLink{sga3:iii:corollary:4-29}{4.29} 和
\SGARefLink{sga3:iii:proposition:4-36}{4.36}).

% Source page: Exp3 local p. 68; combined p. 178; re-edition p. 168;
% original side-pages 148--150.
\medskip
\noindent\SGAInlineRefTarget{sga3:iii:paragraph:4-24}\textbf{4.24. ---}
沿用 \SGARefLink{sga3:iii:theorem:4-21}{4.21} 的记号.
设 \(Y\) 是 \(X\) 的一个平坦子群, 其约化为 \(Y_{\mathscr J}\).
回顾在 \SGARefLink{sga3:iii:remark:0-5}{0.5} 中引入的函子 \(L_0^X\)
(相应地为 \(L_0^Y\)): 对可变的 \(S_0\)-概形 \(T\), 它由等式
\[
 \operatorname{Hom}_{S_0}(T,L_0^X)
 =\operatorname{Hom}_{\mathscr O_T}
 \left(
  \omega^1_{X_0/S_0}\otimes_{\mathscr O_{S_0}}\mathscr O_T,
  \mathscr J\otimes_{\mathscr O_{S_0}}\mathscr O_T
 \right),
\]
定义; 将 \(X\) 换成 \(Y\) 时也类似. 同时还引入了函子
\[
 L'_X=\prod_{S_0/S}L_0^X.
\]
于是有:

\medskip
\noindent\SGAInlineRefTarget{sga3:iii:lemma:4-25}\textbf{引理 4.25.}
\textit{存在典范的 \(Y_0\)-\(\mathscr O_{S_0}\)-模正合序列}
\[
 \tag{+}\label{sga3:iii:lemma:4-25:equation:plus}
 \mathbf n_{Y_0/X_0}
 \xrightarrow{d}
 \omega^1_{X_0/S_0}
 \longrightarrow
 \omega^1_{Y_0/S_0}
 \longrightarrow 0
\]
\textit{具有下述性质:}
\begin{enumerate}
\item[\SGAInlineRefTarget{sga3:iii:lemma:4-25:item:i}(i)] \textit{拉回到 \(Y_0\) 后,
  \(d\) 给出 \SGARefLink{sga3:iii:proposition:4-8:item:iii}{4.8(iii)} 中的态射
  \(\overline D_0\).}
\item[\SGAInlineRefTarget{sga3:iii:lemma:4-25:item:ii}(ii)] \textit{若 \(X_0\) 和 \(Y_0\)
  在 \(S_0\) 上平滑, 则 \(d\) 为单射. 此时两个微分模均为有限型局部自由模,
  因此 \(\mathbf n_{Y_0/X_0}\) 也是有限型局部自由模, 且这一序列局部分裂.}
\end{enumerate}

\textit{证明.}\footnote{编者注: 这一证明已作扩充, 以计入第 I 篇第 6.8 节增补的内容.}
以 \(\pi_0:Y_0\to S_0\) 表示结构态射.
由 SGA 1 II 的公式 (4.3) (另见 EGA IV\textsubscript{4}, 16.4.21),
存在典范的 \(\mathscr O_{Y_0}\)-模正合序列
\[
 \tag{\(\dagger\)}\label{sga3:iii:lemma:4-25:equation:dagger}
 \mathscr N_{Y_0/X_0}
 \xrightarrow{\overline D_0}
 \Omega^1_{X_0/S_0}\otimes_{\mathscr O_{X_0}}\mathscr O_{Y_0}
 \longrightarrow
 \Omega^1_{Y_0/S_0}
 \longrightarrow 0.
\]
由于这一序列由 \((Y_0\mathbin{\times_S}Y_0)\)-等变模及等变态射组成,
它经 \(\varepsilon_0^*\) 的拉回 \((+)\) 是一个
\(Y_0\)-\(\mathscr O_{S_0}\)-模正合序列,
而 \((\dagger)\) 是 \((+)\) 经 \(\pi_0^*\) 的拉回
(参见\SGARefLink{sga3:i:subsection:6-8}{第 I 篇第 6.8 节}).
这就证明了 \SGARefLink{sga3:iii:lemma:4-25:item:i}{(i)}.

再设 \(X_0\) 和 \(Y_0\) 在 \(S_0\) 上平滑.
则由 SGA 1 II, 4.10 (另见 EGA IV\textsubscript{4}, 17.2.3(i) 和 17.2.5),
\(\overline D_0\) 为单射, 且序列 \((\dagger)\) 由有限型局部自由
\(\mathscr O_{Y_0}\)-模组成 (因此局部分裂).
由范畴等价 I, \SGARefLink{sga3:i:proposition:6-8-1}{6.8.1},
\(d\) 也为单射, 因此序列 \((+)\) 具有所述性质.

% Source page: Exp3 local p. 69; combined p. 179; re-edition p. 169;
% original side-pages 150--151.
\medskip
\noindent\SGAInlineRefTarget{sga3:iii:paragraph:4-26}\textbf{4.26. ---}\footnote{编者注:
已扩充原文, 以明确指出这是一个 \(Y_0\)-\(\mathscr O_{S_0}\)-模正合序列.}
对每个 \(S_0\)-概形 \(f:T\to S_0\), 序列 \((+)\) 给出
\(Y_0(T)\)-\(\mathscr O(T)\)-模正合序列
\[
\begin{aligned}
0&\longrightarrow
 \operatorname{Hom}_{\mathscr O_T}
 \left(f^*\omega^1_{Y_0/S_0},f^*\mathscr J\right)
 \longrightarrow
 \operatorname{Hom}_{\mathscr O_T}
 \left(f^*\omega^1_{X_0/S_0},f^*\mathscr J\right)\\
&\hspace{10em}\longrightarrow
 \operatorname{Hom}_{\mathscr O_T}
 \left(f^*\mathbf n_{Y_0/X_0},f^*\mathscr J\right).
\end{aligned}
\]
因此存在 \(Y_0\)-\(\mathscr O_{S_0}\)-模正合序列
\[
 \tag{4.26.1}\label{sga3:iii:paragraph:4-26:equation:4-26-1}
 0\longrightarrow L_0^Y\longrightarrow L_0^X
 \xrightarrow{d}N_0.
\]
它给出一个交换群复形的正合序列
\[
 0\longrightarrow C^*(Y_0,L_0^Y)
 \longrightarrow C^*(Y_0,L_0^X)
 \xrightarrow{d^*}C^*(Y_0,N_0),
\]
特别地, 得到以下各行正合的交换图:

% Raster stand-in for Loop 1; replace by a native TeX diagram in Loop 2.
% Authority: Exp3 local p. 69; combined p. 179; re-edition p. 169;
% original side-page 150.
\begin{center}
 \SGAThreeNativeInput[0.88\textwidth]
   {figures/exp3/diagram_060_exp3_p069_exact_complex_rows.tex}
\end{center}

注意到 \(C^0(Y_0,L_0^Y)\) (相应地为 \(C^0(Y_0,L_0^X)\)) 正是
\[
 \operatorname{Hom}_{S_0}(S_0,L_0^Y)
 =\operatorname{Hom}_{S}(S,L'_Y)
\]
(相应地是将 \(Y\) 换成 \(X\) 所得的类似群);
由 \SGARefLink{sga3:iii:corollary:0-9}{0.9}, 这是 \(Y\) (相应地为 \(X\))
在 \(S\) 上诱导 \(X_{\mathscr J}\) 的单位截面的那些截面组成的群.
此外, \(d^1\) 恰好是 \SGARefLink{sga3:iii:proposition:4-8:item:iii}{4.8(iii)}
中的态射 \(v_{i_{Y_0}}\), 其中 \(i_{Y_0}:Y_0\to X_0\) 为典范浸入.
\footnote{编者注: 这由 \(d:\mathbf n_{Y_0/X_0}\to\Omega^1_{X_0/S_0}\)
的定义 (参见 \SGARefLink{sga3:iii:lemma:4-25}{4.25}) 及 \(v_{i_{Y_0}}\)
的定义 (参见 \SGARefLink{sga3:iii:proposition:4-8}{4.8}) 得到.}

\medskip
\noindent\SGAInlineRefTarget{sga3:iii:lemma:4-27}\textbf{引理 4.27.}
\textit{设 \(S\), \(\mathscr I\), \(\mathscr J\) 和 \(X\) 满足
\SGARefLink{sga3:iii:theorem:4-21}{4.21} 的条件, 并设 \(Y\) 是 \(X\) 的子群,
它在 \(S\) 上平坦且提升 \(Y_{\mathscr J}\).
以 \(i:Y\hookrightarrow X\) 表示典范浸入.}
\footnote{编者注: 下面的 \SGARefLink{sga3:iii:lemma:4-27:item:i}{(i)} 为增补内容;
它将在 \SGARefLink{sga3:iii:proposition:4-35-1}{4.35.1},
继而在 \SGARefLink{sga3:iii:paragraph:4-38:item:4}{4.38(4)} 及
\SGARefLink{sga3:iii:paragraph:4-38:item:5}{(5)} 中用到.}
\begin{enumerate}
\item[\SGAInlineRefTarget{sga3:iii:lemma:4-27:item:i}(i)] \textit{设 \(i':Y\to X\)
  是提升 \(i_0\) 的 \(S\)-概形态射 (因而 \(i'\) 也是浸入),
  置 \(Y'=i'(Y)\), 并以 \(d(i,i')\) 表示使 \(i'=d(i,i')\cdot i\)
  成立的 \(C^1(Y_0,L_0^X)\) 的元素 (参见 \SGARefLink{sga3:iii:proposition:1-2-4}{1.2.4}).
  则偏差 \(d(Y,Y')\in C^1(Y_0,N_0)\)
  (参见 \SGARefLink{sga3:iii:remark:4-5-1}{4.5.1}) 由下式给出:}
  \[
   d(Y,Y')=d^1\bigl(d(i,i')\bigr).
  \]
\item[\SGAInlineRefTarget{sga3:iii:lemma:4-27:item:ii}(ii)] \textit{设 \(x\in C^0(Y_0,L_0^X)\)
  是 \(X\) 在 \(S\) 上的一个截面, 它在 \(S_{\mathscr J}\) 上诱导
  \(X_{\mathscr J}\) 的单位截面. 则偏差
  \(d(Y,\operatorname{Int}(x)Y)\in C^1(Y_0,N_0)\)
  (参见 \SGARefLink{sga3:iii:remark:4-5-1}{4.5.1}) 由下式给出:}
  \[
   -d\bigl(Y,\operatorname{Int}(x)Y\bigr)
   =d^1\partial x=\partial d^0x.
  \]
\end{enumerate}

% Source page: Exp3 local p. 70; combined p. 180; re-edition p. 170;
% original side-page 152.
事实上, \(Y'\) 是 \(Y\) 在复合态射
\footnote{编者注: 回顾 \(L'_X=\prod_{S_0/S}L_0^X\).}
\[
 Y\xrightarrow{(d(i,i'),i)}L'_X\mathbin{\times_S}X
 \longrightarrow X,
\]
下的像, 此复合态射在 \SGARefLink{sga3:iii:proposition:4-8:item:iii}{4.8(iii)} 中记为
\(d(i,i')\cdot i\). 因此, 由该处的结果及等式 \(v_{i_0}=d^1\), 有
\[
 \begin{aligned}
 &c\bigl(X,Y',d(i,i')\cdot i\bigr)-c(X,Y',i)\\
 &\qquad=v_{i_0}\bigl(d(i,i')\bigr)
 =d^1\bigl(d(i,i')\bigr).
 \end{aligned}
\]
但按定义, \(d(i,i')\cdot i=i'\) 经过 \(Y'\) 分解, 所以第一项为零.
此外, 由 \SGARefLink{sga3:iii:proposition:4-8:item:ii}{4.8(ii)}, 有
\[
 c(X,Y',i)=d(Y',Y)=-d(Y,Y').
\]
因而 \(d(Y,Y')=d^1(d(i,i'))\), 这就证明了
\SGARefLink{sga3:iii:lemma:4-27:item:i}{(i)}.

现设 \(x\) 如 \SGARefLink{sga3:iii:lemma:4-27:item:ii}{(ii)} 所述. 由公式
\[
 xyx^{-1}=xyx^{-1}y^{-1}y
 =(x-\operatorname{Ad}(y)x)y=(-\partial x)(y)\cdot y,
\]
可知 \(Y'\) 是 \(Y\) 在浸入 \(i'=(-\partial x)\cdot i_Y\) 下的像.
于是由 \SGARefLink{sga3:iii:lemma:4-27:item:i}{(i)} 得到
\[
 -d\bigl(Y,\operatorname{Int}(x)Y\bigr)
 =d^1\partial x=\partial d^0x.
\]

\medskip
\noindent\SGAInlineRefTarget{sga3:iii:corollary:4-28}\textbf{推论 4.28.}
\textit{设 \(Y\) 和 \(Y'\) 是 \(X\) 的两个子群,
它们在 \(S\) 上平坦且提升 \(Y_{\mathscr J}\).
它们由 \(X\) 在 \(S\) 上的一个诱导 \(X_{\mathscr J}\) 的单位截面的截面共轭,
当且仅当}
\[
 d(Y,Y')\in
 \partial d^0 C^0(Y_0,L_0^X)
 \subset \partial C^0(Y_0,N_0)
 =B^1(Y_0,N_0).
\]
\noindent\emph{底本说明.} 法文重版本此处印为 \(G\);
本节中的环境群始终是 \(X\).

\medskip
\noindent\SGAInlineRefTarget{sga3:iii:corollary:4-29}\textbf{推论 4.29.}
\textit{若 \(d^0\) 为满射, 则上述 \(Y\) 和 \(Y'\) 由 \(X\) 在 \(S\) 上的
一个诱导 \(X_{\mathscr J}\) 的单位截面的截面共轭, 当且仅当
\(d(Y,Y')\in B^1(Y_0,N_0)\).}

\medskip
\noindent\SGAInlineRefTarget{sga3:iii:corollary:4-30}\textbf{推论 4.30.}
\textit{设 \(Y\) 如 \SGARefLink{sga3:iii:lemma:4-27}{4.27} 所述.
用 \(X\) 在 \(S\) 上诱导 \(X_{\mathscr J}\) 的单位截面的截面对 \(Y\) 作共轭,
所得共轭子群的集合与下列集合同构:}
\[
 d^1\partial C^0(Y_0,L_0^X)
 =C^0(Y_0,L_0^X)\big/\ker(d^1\partial).
\]

现在注意到, \(C^0(Y_0,L_0^X)/\ker(d^1\partial)\) 的计算,
仅涉及 \SGARefLink{sga3:iii:paragraph:4-26}{4.26} 的交换图中左边的方块.
特别地, 也可以在任何具有相同左方块的同类图中计算它.
考虑 \(S_0\) 上的函子 \(L_0^X/L_0^Y\), 其定义为
\[
 \operatorname{Hom}_{S_0}(T,L_0^X/L_0^Y)
 =\operatorname{Hom}_{S_0}(T,L_0^X)
  \big/\operatorname{Hom}_{S_0}(T,L_0^Y).
\]
存在交换图

% Raster stand-in for Loop 1; replace by a native TeX diagram in Loop 2.
% Authority: Exp3 local p. 70; combined p. 180; re-edition p. 170;
% original side-page 152.
\begin{center}
 \SGAThreeNativeInput[0.91\textwidth]
   {figures/exp3/diagram_061_exp3_p070_quotient_exact_complex_rows.tex}
\end{center}

由前面的观察得到:

% Continue with Corollary 4.31 on Exp3 local p. 71 / combined p. 181 /
% re-edition p. 171 / original side-page 153.
